\chapter{Fixed Points}

\section{What "Nothing Happens" Actually Means}

Chapter 15 closed with a brief, deliberately informal image: a courtyard region so settled, so densely supported by accumulated residue and so thoroughly resolved in appearance, that whatever continues to be proposed there tends to commit back to what was already true. It is worth pausing on that image before formalizing it, because "nothing happens" turns out to describe two quite different situations, and this book's framework needs to keep them apart.

The first situation is genuine inertness. No \(\Delta\psi\) is proposed at this region at all; nothing presses on it, nothing forces, and the field simply holds its last committed value because nothing has asked it to do otherwise. A stretch of bare, undisturbed paving, far from any doorway or bench, well outside the reach of any cue strong enough to activate a proposal, is stable in this sense: stable by default, through absence of pressure rather than through any active work.

The second situation looks identical from a distance and is, in an important sense, the opposite. Consider a fountain at the courtyard's center. Water moves through it continuously; \(\Delta\psi\) is being proposed there at every update, not occasionally but constantly, since flow is precisely what a fountain's vector field v is supposed to represent. And yet the fountain, observed a minute apart, a day apart, a year apart, is recognizably the same fountain, its form unchanged even though the specific water occupying it has been replaced many times over. This is stability achieved through continuous, active maintenance rather than through absence of proposal, and this chapter's task is to give this second kind of stability, the more interesting and by far the more common kind for anything this book would want to call an object, a precise formal description.

\section{The Fixed Point Condition}

Both situations satisfy a single equation, once the write cycle from Chapters 12 through 15 is stated in its settled form. A field state \(M^*\) is a fixed point of the update cycle when projecting and committing whatever it proposes for itself returns the same state:

\[ M^* = \Pi_C(M^* + \Delta\psi). \]

This is the same condition the supporting technical material behind this book states in two closely related forms, \(M^*\) = Update(\(M^*\)) in the architecture specification's language and \(M^*\) = \(\Pi_C\)(\(M^*\) + \(\Delta\psi\)) in the dynamics module's, and this chapter adopts the second form because it makes explicit what the first leaves implicit: reaching a fixed point is not a matter of nothing being proposed, but of whatever is proposed surviving projection and commit unchanged.

Read this way, the equation does not distinguish the bare paving from the fountain. Both satisfy it. What distinguishes them is what \(\Delta\psi\) actually is at each: at the bare paving, \(\Delta\psi\) is zero or close to it, and the equation holds trivially, because there is nothing to project or commit in the first place. At the fountain, \(\Delta\psi\) is substantial and constantly regenerated, and the equation holds nontrivially, because the specific flow proposed at each update, once projected against the field's constraints and committed, happens to reproduce the fountain's form rather than drifting away from it. The next section makes this distinction explicit, because collapsing it back together would undo the very thing this chapter exists to separate.

\section{Two Kinds of Fixed Point}

Call a fixed point trivial when \(\Delta\psi\) = 0 at every point it covers: nothing is proposed, so nothing needs to survive projection, and the state persists purely by default. Call a fixed point dynamic when \(\Delta\psi\) is genuinely nonzero, sometimes substantially so, but the cycle of proposal, projection, and commit reliably returns the same state regardless. The fountain is a dynamic fixed point. So, in fact, is the stone wall this book opened with, once its full history is taken into account: weathering proposes small changes to its surface continuously, residue accumulates, and the wall's form nonetheless persists because whatever is proposed for it projects back onto essentially the same admissible configuration, update after update.

This distinction matters because dynamic fixed points are doing something trivial fixed points are not. A trivial fixed point persists because it costs nothing to maintain; there is, in a real sense, nothing actively holding it in place beyond the absence of disturbance. A dynamic fixed point persists despite continuous pressure to become something else, and its stability is therefore a genuine achievement of the field's own dynamics rather than an accident of neglect. Most of what this book, and its reader, would naturally call an object, a wall, a bench, a fountain, a courtyard as a whole, turns out on inspection to be a dynamic fixed point rather than a trivial one, however static it may appear to casual observation.

\section{Identity as a Fixed Point, Not a Label}

This closes a loop left open since Chapter 6. Section 6.6 argued that identity, in a field-based ontology, cannot be a label attached to a region, and had to instead be something the field's own dynamics actively maintained: a coherent, self-reinforcing pattern that persists because the field's evolution keeps reproducing it, not because something outside the field decided to keep calling it by the same name. That description was necessarily informal at the time, since neither \(\Delta\psi\), projection, nor commit had yet been introduced. It can now be stated exactly. An object, in this framework, is a dynamic fixed point of the write cycle: a region where \(M^*\) = \(\Pi_C\)(\(M^*\) + \(\Delta\psi\)) holds nontrivially, where proposal continues to arrive and continues to be absorbed back into the same admissible configuration rather than drifting away from it. The wall is the wall not because a label says so, but because it is, quite literally, a solution to this chapter's equation, reproduced update after update against a continuous stream of proposals that could, in principle, have taken it somewhere else and consistently did not.

\section{Stability and Its Failure}

Not every fixed point is equally robust, and it is worth flagging the distinction here even though its full treatment belongs to Part IV. A fixed point is stable, informally, when a small perturbation, a proposal slightly outside what the region has settled into, gets absorbed back toward \(M^*\) on the next cycle rather than growing into something else. A fixed point is unstable when the opposite happens: a small perturbation, rather than being corrected, is amplified by the next round of proposal and projection, carrying the region away from \(M^*\) rather than back toward it. A well-built wall is a stable fixed point in this sense; a precarious stack of loose stones, satisfying the fixed-point equation for as long as nothing disturbs it, may be an unstable one, technically a solution to \(M^*\) = \(\Pi_C\)(\(M^*\) + \(\Delta\psi\)) while remaining one small proposal away from becoming a very different solution entirely. This book will not develop the mathematics of stability further here; Part IV's treatment of variational dynamics and constraint manifolds is where the distinction between a fixed point that recovers from disturbance and one that merely has not yet been disturbed receives its proper formal footing.

\section{Looking Ahead}

This chapter has treated the field's update cycle as though it operated on one uniform clock, proposal, projection, and commit happening together, at the same rate, everywhere. This was a simplification, and Part I's original architecture specification already anticipated why it could not hold in general: appearance and flow tend to change quickly, while coherence, structure, and identity tend to change slowly, and a region can be a fixed point with respect to one of these timescales while still actively settling with respect to another. A fountain's water is never at rest, yet the fountain's overall form was already treated in this chapter as fixed; those two claims are only compatible once fast and slow dynamics are allowed to operate on separate clocks rather than one. That separation, and what it means for a single region to be simultaneously settled and still-evolving, is the subject of Chapter 17.
